% ===================================================================
%  FIN DU LIVRET — feuillets 109 a 116 : TABELO DIL KONTENAJO, pages
%  d'annonce de l'editeur, acheve d'imprimer.
%
%  Aucun de ces feuillets ne porte de folio imprime (verifie sur le
%  scan : ni haut ni bas de page). Les feuillets 114 et 115 sont des
%  pages blanches (0,1 % d'encre, grain de papier seul).
%
%  Feuillets 110 a 113 et 116 : pages d'annonce de l'editeur Delmas,
%  imprimees en FRANCAIS (non traduites), sur une justification plus
%  large que le corps du livre (~86 mm contre 79,4 mm) : d'ou des
%  corps PROVISOIRES plus petits que ceux du texte courant, choisis
%  pour que chaque ligne relevee tienne dans \VUtexteLargeur sans
%  Overfull. A reprendre par outils/apparat.py comme le reste de
%  l'apparat.
% ===================================================================

% --- Feuillet 109 (TABELO DIL KONTENAJO + acheve d'imprimer, non
%     foliote) -------------------------------------------------------
\begin{VUpage}[109]{}
%%K t99-tit-1 sub
\VUsaut{3.0mm}
\VUtitre{15.0pt}{30}{TABELO DIL KONTENAJO}
\VUsaut{2.4mm}
\VUfilet{20mm}
\VUsaut{5.0mm}
{\raggedleft\textit{Pages}\par}
\VUsaut{1.0mm}

{\fontsize{10.2pt}{12.2pt}\selectfont
%%K t99-01-1 p
\noindent \textsc{Avizo} \dotfill 3

%%K t99-02-1 p
\noindent \VUgras{Tabelo 1.} --- La skolo. --- La liceo. --- La doco\cc
chambro \dotfill 5

%%K t99-03-1 p
\noindent \VUgras{Tabelo 2.} --- La homala korpo. --- L'amuztempo.\nl
--- La ludi \dotfill 11

%%K t99-04-1 p
\noindent \VUgras{Tabelo 3.} --- La puereso. --- (Bapto-ceremonio en\nl
vilajo) \dotfill 17

%%K t99-05-1 p
\noindent La yuneso (publika festo) \dotfill 20

%%K t99-06-1 p
\noindent \VUgras{Tabelo 4.} --- La evo-matureso (la mariaj-festino) \dotfill 24

%%K t99-07-1 p
\noindent La oldeso (aniversario dil avulo) \dotfill 27

%%K t99-08-1 p
\noindent \VUgras{Tabelo 5.} --- La domo e lua konstrukto \dotfill 30

%%K t99-09-1 p
\noindent \VUgras{Tabelo 6.} --- La internajo dil domo. --- La moblaro.\nl
La nutrado. --- La lumizado \dotfill 37

%%K t99-10-1 p
\noindent \VUgras{Tabelo 7.} --- La vilajo (vintre) \dotfill 47

%%K t99-11-1 p
\noindent La rustika domo (printempe) \dotfill 51

%%K t99-12-1 p
\noindent \VUgras{Tabelo 8.} --- La rekolto (somere) \dotfill 55

%%K t99-13-1 p
\noindent La vit-rekolto (autune) \dotfill 58

%%K t99-14-1 p
\noindent \VUgras{Tabelo 9.} --- La monto, la baln-urbo \dotfill 61

%%K t99-15-1 p
\noindent La foresto, la chasado \dotfill 64

%%K t99-16-1 p
\noindent \VUgras{Tabelo 10.} --- La maro, la portuo \dotfill 68

%%K t99-17-1 p
\noindent \VUgras{Tabelo 11.} --- La urbo e lua monumenti. --- L'in\cc
cendio \dotfill 75

%%K t99-18-1 p
\noindent \VUgras{Tabelo 12.} --- La staciono. --- La fervoyi \dotfill 80

%%K t99-19-1 p
\noindent \VUgras{Tabelo 13.} --- La hotelo, la restorerio, la kafeerio \dotfill 85

%%K t99-20-1 p
\noindent \VUgras{Tabelo 14.} --- La strado, la komercisti \dotfill 89

%%K t99-21-1 p
\noindent \VUgras{Tabelo 15.} --- La merkato, la manjebli \dotfill 96

%%K t99-22-1 p
\noindent \VUgras{Tabelo 16.} --- La granda magazino, la bazaro, la\nl
ludili \dotfill 101}

%%K t99-apar-1 apar
\VUsaut{8.0mm}
\VUfilet{20mm}
\VUsaut{9.0mm}
\VUcentre{8.0pt}{50}{\textsc{Imprimerie des Tableaux Auxiliaires Delmas}}
\VUsaut{1.6mm}
\VUcentre{8.0pt}{0}{6, Place Saint-Christoly, 6}
\VUsaut{1.2mm}
\VUcentre{8.0pt}{50}{\textsc{Bordeaux-France}}
\VUsaut{1.6mm}
\VUfilet{5mm}
\VUsaut{1.6mm}
\VUcentre{8.0pt}{0}{1926}
\end{VUpage}

% --- Feuillet 110 (annonce Delmas, non foliote) ---------------------
%  Publicite pour les « T. A. D. », en francais, sur mesure plus
%  large que le corps du livre.
\begin{VUpage}[110]{}
%%K t99-apar-2 apar
% PLANCHE D'ANNONCE. Ses titres et ses courtes lignes d'apparat se
% composent en \VUtitre et \VUcentre ; ses paragraphes sont du texte
% courant qui remplit reellement la justification, et gardent donc
% leurs coupures \nl. Aucune ligne de cette page n'est une ligne
% d'affiche etiree a tort : rien n'appelle \VUlineo ici.
\VUsaut{2.0mm}
\VUcentre{9.0pt}{0}{\textit{Pour l'Enseignement pratique des Langues vivantes}}
\VUsaut{3.0mm}
\VUtitre{15.0pt}{60}{les " T. A. D. "}
\VUsaut{2.4mm}
\VUtitre{11.5pt}{20}{Tableaux Auxiliaires Delmas}
\VUsaut{3.0mm}
\VUcentre{8.5pt}{0}{\textit{sont employés dans plus de :}}
\VUsaut{1.6mm}
\VUcentre{9.5pt}{20}{\textit{\VUgras{8,000 Ecoles du Monde entier}}}
\VUsaut{3.0mm}

{\fontsize{8.5pt}{10.2pt}\selectfont
\noindent \VUgras{Des milliers de références, que nous tenons à votre}\nl
\VUgras{disposition, montrent que la « Collection des}\nl
\VUgras{T. A. D. » est l'auxiliaire indispensable de toute}\nl
\VUgras{méthode d'enseignement des Langues vivantes, quel}\nl
\VUgras{que soit le Programme Officiel imposé.}}

\VUsaut{5.0mm}
\VUtitre{10.5pt}{20}{Pour le Professsseur en Classe}
\VUsaut{2.6mm}
\VUtitre{13.0pt}{20}{16 Tableaux Muraux en Couleurs}
\VUsaut{2.4mm}
\VUcentre{9.5pt}{0}{\VUgras{de dimensions de 90 $\times$ 120 centimètres}}
\VUsaut{3.0mm}

{\fontsize{8.5pt}{10.2pt}\selectfont
\noindent permettent au maître :

\VUblancAlinea
de \VUgras{fixer l'attention de tous ses élèves} sur un sujet\nl
de cours oral, l'image représentée sur le Tableau;

\VUblancAlinea
d'énoncer \VUgras{les mots} avec leur prononciation cor\cc
recte en montrant \VUgras{l'objet} qu'ils représentent;

\VUblancAlinea
Les Elèves \VUgras{entendent} et \VUgras{voient} et \VUgras{apprennent}\nl
ainsi, agréablement et avec beaucoup plus de facilité,\nl
le \VUgras{vocabulaire} et la \VUgras{prononciation} et insensiblement,\nl
automatiquement l'emploi des \VUgras{mots} dans la phrase,\nl
\VUgras{l'emploi des verbes.}}

\end{VUpage}

% --- Feuillet 111 (gravure + annonce, non foliote) ------------------
%  Le haut de page porte une gravure (reduction en noir du tableau
%  mural n° 12) : non reproduite. Hauteur estimee sur le rapport
%  d'encre par ligne du scan brut (bloc dense jusque vers 46 % de la
%  hauteur de page, puis blanc avant la legende).
\begin{VUpage}[111]{}
%%K t99-apar-3 apar
% LA GRAVURE Y EST MAINTENANT. Elle etait remplacee par un blanc de
% 85 mm et un commentaire : la page annoncait « Reduction du Tableau
% Mural n^o 12 » au-dessous d'un vide. L'image est prelevee sur le
% feuillet meme (outils/imaji.py) et posee a sa largeur relevee.
\VUsaut{2.0mm}
\VUimajo{tabelo-12.png}{1.0}
\VUsaut{1.6mm}
\VUcentre{8.0pt}{0}{Réduction du Tableau Mural n\textsuperscript{o} 12 « LES VOYAGES »}
\VUsaut{1.0mm}
\VUcentre{8.0pt}{0}{Dimensions du Tableau Mural : 90 centimètres sur 120 centimètres.}
\VUsaut{4.0mm}

{\fontsize{8.5pt}{10.2pt}\selectfont
\noindent \textit{Le Maître peut, suivant sa méthode et ses idées per-}\nl
\textit{sonnelles, suivant les progrès des élèves, développer}\nl
\textit{comme il le juge utile les sujets fixés sur les 16 Ta-}\nl
\textit{bleaux qui représentent :}

\VUblancAlinea
N\textsuperscript{o} 1. --- L'Ecole, le Lycée, la Classe, les Nombres.

\VUblancAlinea
2. --- La Récréation, les Jeux, le Corps humain.

\VUblancAlinea
3. --- Enfance et Jeunesse : \textit{Le Temps, l'Age, la}\nl
\textit{Famille.}

\VUblancAlinea
4. --- Age mûr et Vieillesse : \textit{Le Vêtement, la}\nl
\textit{Santé.}

\VUblancAlinea
5. --- La Maison et sa Construction, \textit{les Corps de}\nl
\textit{métiers, les Ouvriers, les Outils.}

\VUblancAlinea
6. --- L'intérieur de la Maison, \textit{les Pièces et Meu-}\nl
\textit{bles, la Nourriture, la Vie de Famille.}}
\end{VUpage}

% --- Feuillet 112 (suite annonce, non foliote) -----------------------
\begin{VUpage}[112]{}
%%K t99-apar-4 apar
{\fontsize{8.5pt}{10.2pt}\selectfont
\noindent 7. --- Le Village en Hiver, la Maison rustique au\nl
\textit{Printemps,} les Métiers divers.

\VUblancAlinea
11. --- La Ville, \textit{Vue d'ensemble, Rues, Moyens de}\nl
\textit{Communication, Monuments, etc.}

\VUblancAlinea
12. --- Les Voyages, \textit{le Chemin de fer, les Ba-}\nl
\textit{teaux, etc.}

\VUblancAlinea
14. --- La Rue et les Commerçants, \textit{etc.}

\VUblancAlinea
15. --- Le Marché, les Comestibles.

\VUblancAlinea
16. --- Un Grand Magasin, \textit{les Objets de toutes}\nl
\textit{sortes.}

\VUblancAlinea
\textit{Une même collection des Tableaux Muraux peut ser-}\nl
\textit{vir pour tous les Professeurs de Langues Vivantes}\nl
\textit{d'un même Collège.}}

\VUsaut{6.0mm}
\VUtitre{10.5pt}{20}{Pour l'Elève, à l'Etude ou chez lui}
\VUsaut{3.0mm}
\VUtitre{14.5pt}{20}{Deux Cahiers Portatifs}
\VUsaut{2.2mm}
\VUcentre{9.5pt}{0}{\VUgras{de dimensions 24 $\times$ 32 centimètres}}
\VUsaut{3.0mm}

{\fontsize{8.5pt}{10.2pt}\selectfont
\noindent contenant:

\VUblancAlinea
la \VUgras{reproduction réduite} et en noir des 16 Tableaux\nl
Muraux;

\VUblancAlinea
les \VUgras{vocabulaires ou listes de mots} correspondant\nl
aux Tableaux avec les chiffres correspondant aux objets\nl
représentés,

\VUblancAlinea
permettent à l'Elève qui a appris en classe les mots,\nl
leur prononciation, d'en acquérir \VUgras{l'orthographe cor}\cc
\VUgras{recte} et de faire, d'après ces Tableaux et Vocabulai\cc
res, \VUgras{les devoirs} indiqués par le Professeur.

\VUblancAlinea
\VUgras{\textit{Ces Vocabulaires sont édités en Français, Anglais,}}\nl
\VUgras{\textit{Allemand, Espagnol, Italien, Russe ou Flamand.}}}
\end{VUpage}

% --- Feuillet 113 (suite annonce, non foliote) -----------------------
\begin{VUpage}[113]{}
%%K t99-apar-5 apar
% PLANCHE D'ANNONCE. Comme le feuillet 110, ses titres et lignes
% d'apparat se composent en \VUtitre et \VUcentre ; ses paragraphes
% remplissent reellement la justification et gardent leurs \nl. Seule
% particularite : « Mon { Premier / Deuxieme / Troisieme / Quatrieme }
% Livre de Francais » porte une vraie accolade, composee plus bas en
% array + \left\{ ... \right\}, non decrite en prose.
\VUsaut{2.0mm}
\VUtitre{10.5pt}{20}{Pour le Maître et l'Elève}
\VUsaut{3.0mm}
\VUtitre{12.5pt}{20}{Le " Livret explicatif des T. A. D "}
\VUsaut{3.0mm}

{\fontsize{8.5pt}{10.2pt}\selectfont
\noindent \textit{contenant la} \VUgras{\textit{description complète}} \textit{des 16 Tableaux}\nl
\textit{avec emploi des mots, des verbes, etc., est :}

\VUblancAlinea
\textit{avant le Cours, le} \VUgras{\textit{Guide du Maître;}}

\VUblancAlinea
\textit{après le Cours, le} \VUgras{\textit{Livre de Lecture de l'Elève.}}

\VUblancAlinea
\textit{Comme pour les Vocabulaires, un Livret explicatif}\nl
\textit{est édité dans chacune des Langues ci-dessus indi-}\nl
\textit{quées.}}

\VUsaut{6.0mm}

{\fontsize{8.5pt}{10.2pt}\selectfont
\noindent \VUgras{Pour l'enseignement de la Langue Française, la}\nl
\VUgras{« Collection des T. A. D. » a édité un :}}

\VUsaut{3.0mm}
\VUtitre{7.2pt}{0}{Cours Pratique d'Enseignement de la Langue Française}
\VUsaut{1.6mm}
\VUtitre{11.5pt}{20}{aux Étrangers}
\VUsaut{2.4mm}
{\fontsize{6.5pt}{7.8pt}\selectfont\VUcentre{6.5pt}{0}{\VUgras{par la Méthode Directe et d'après les Tabl aux Auxi'iaires Delmas.}}}
\VUsaut{3.0mm}

{\fontsize{8.5pt}{10.2pt}\selectfont
\noindent \VUgras{Ce Cours est complet (vocabulaire et grammaire),}\nl
\VUgras{en quatre livres intitulés :}}

\VUsaut{3.0mm}
% « Mon { Premier / Deuxième / Troisième / Quatrième } Livre de Français » :
% une accolade rassemble les quatre titres entre « Mon » et « Livre de
% Français », comme au fac-simile -- ce n'est pas une ligne de plus a
% decrire, c'est une vraie accolade a composer. L'environnement array,
% avec \left\{ ... \right\}, la dessine a la hauteur des quatre lignes.
% \lineskiplimit=-1000pt (impose par VUpage, cf. preambule.tex § 2) force
% PARTOUT l'interligne fixe de la page : a l'interieur d'un array, cela
% ecrasait l'espacement des rangs (les quatre mots s'imprimaient l'un
% sur l'autre) -- corrige en restaurant la grille normale de LaTeX,
% localement, pour la duree du array. Meme cause, deuxieme effet : la
% ligne qui SUIT le bloc heritait d'un \prevdepth enorme (le bloc fait
% quatre rangs de haut) et le meme calcul rendait un blanc NEGATIF --
% « par M. E. Rochelle, » s'imprimait par-dessus « Quatrième ». On sort
% donc aussi ce bloc de la grille par un \nointerlineskip de chaque
% cote, et on lui rend un blanc pose a la main, comme un \VUimajo.
\nointerlineskip
{\lineskiplimit=0pt\lineskip=1pt
\noindent\hbox to \VUtexteLargeur{\hfil
  {\fontsize{9.5pt}{9.5pt}\selectfont\VUgras{Mon}}\hspace{2.5mm}%
  $\left\{\begin{array}{c}
    \mbox{\fontsize{9.5pt}{9.5pt}\selectfont\VUgras{Premier}}\\[3.2pt]
    \mbox{\fontsize{9.5pt}{9.5pt}\selectfont\VUgras{Deuxième}}\\[3.2pt]
    \mbox{\fontsize{9.5pt}{9.5pt}\selectfont\VUgras{Troisième}}\\[3.2pt]
    \mbox{\fontsize{9.5pt}{9.5pt}\selectfont\VUgras{Quatrième}}
  \end{array}\right\}$\hspace{2.5mm}%
  {\fontsize{9.5pt}{9.5pt}\selectfont\VUgras{Livre de Français}}%
\hfil}\par}
\nointerlineskip
\VUsaut{2.6mm}
\VUcentre{8.5pt}{0}{par M. E. Rochelle,}
\VUsaut{1.2mm}
\VUcentre{8.5pt}{0}{Professeur au Lycée de Bordeaux.}
\VUsaut{2.4mm}
\VUfilet{18mm}
\VUsaut{0.4mm}
\VUfilet{18mm}
\VUsaut{4.0mm}
\VUcentre{9.0pt}{0}{\textit{Demandez aux}}
\VUsaut{2.0mm}
\VUtitre{13.0pt}{20}{Tableaux Auxiliaires Delmas}
\VUsaut{1.6mm}
\VUcentre{9.0pt}{0}{6, Place Saint-Christoly, \VUgras{Bordeaux (France)}}
\VUsaut{3.0mm}

{\fontsize{8.5pt}{10.2pt}\selectfont
\noindent des \VUgras{spécimens, tarifs} et tous renseignements com\cc
plémentaires.

\VUblancAlinea
Ils vous seront adressés gratuitement et sans enga\cc
gement de votre part.}

\end{VUpage}

% --- Feuillets 114-115 (pages blanches, non foliotees) ---------------
\begin{VUpage}[114]{}
% (page vierge)
\end{VUpage}

\begin{VUpage}[115]{}
% (page vierge)
\end{VUpage}

% --- Feuillet 116 (4e de couverture, non foliote) ---------------------
%  Repete le colophon du feuillet 109, en pied de couverture ; le
%  reste de la page est blanc.
\begin{VUpage}[116]{}
%%K t99-apar-6 apar
\VUsaut{90.0mm}
\VUcentre{8.0pt}{50}{\textsc{Imprimerie des Tableaux Auxiliaires Delmas}}
\VUsaut{1.6mm}
\VUcentre{8.0pt}{0}{6, Place Saint-Christoly, 6}
\VUsaut{1.2mm}
\VUcentre{8.0pt}{50}{\textsc{Bordeaux-France}}
\VUsaut{1.6mm}
\VUfilet{5mm}
\VUsaut{1.6mm}
\VUcentre{8.0pt}{0}{1926}
\end{VUpage}
